% =============================================================================
% Conclusion Générale
% =============================================================================
\chapter*{Conclusion Générale}
\addcontentsline{toc}{chapter}{Conclusion Générale}
\markboth{Conclusion Générale}{Conclusion Générale}

\section*{Bilan du projet}

Le stage de fin d'études effectué au sein de l'entreprise DevWise à Nabeul avait
pour objectif principal la mise en place d'un robot de tests automatisés pour les
projets Speed Line et Delivery App. Au terme de ce stage de quatre mois, nous
pouvons dresser un bilan positif des réalisations accomplies.

Le robot de tests automatisés que nous avons conçu et développé couvre
l'intégralité de la pyramide de tests, conformément aux bonnes pratiques de
l'ingénierie logicielle~:

\begin{itemize}[leftmargin=2cm]
  \item \textbf{Tests unitaires} --- 419 tests répartis sur les 12 microservices
    du backend (JUnit~5 + Mockito) et 34 tests pour les deux applications
    Angular frontend (Jasmine + Karma), avec une couverture de code mesurée
    par JaCoCo et Istanbul.

  \item \textbf{Tests d'intégration et API} --- 104 tests Playwright couvrant
    l'ensemble des endpoints REST des 12 microservices, incluant des tests
    de contrats API, des tests de smoke et des tests de WebSocket.

  \item \textbf{Tests end-to-end} --- 66 tests d'interface utilisateur pour le
    panneau d'administration et le tableau de bord partenaire, ainsi que des
    tests d'intégration Flutter avec Patrol pour les applications mobiles
    client et livreur. Des tests de régression visuelle et d'accessibilité
    WCAG~2.1 complètent cette couche.

  \item \textbf{Tests de performance} --- 3 scénarios k6 (charge, stress et
    spike) permettant de valider le comportement du système sous différentes
    conditions de trafic.

  \item \textbf{Tests de sécurité} --- 16 tests vérifiant la résistance aux
    injections SQL, aux attaques XSS, au contournement d'authentification,
    au rate limiting et à l'invalidation des tokens JWT.
\end{itemize}

\section*{Apports du projet}

L'impact du robot de tests sur la qualité des projets Speed Line et Delivery App
se mesure à travers plusieurs indicateurs~:

\begin{table}[H]
\centering
\caption{Comparaison avant/après la mise en place du robot de tests}
\begin{tabular}{|l|c|c|}
\hline
\rowcolor{speedlineBlue!10}
\textbf{Métrique} & \textbf{Avant} & \textbf{Après} \\
\hline
Services avec tests unitaires & 3/12 & 12/12 \\
\hline
Fichiers de test Java & 63 & 84 \\
\hline
Tests unitaires (total) & $\sim$200 & 419+ \\
\hline
Tests API automatisés & 0 & 104 \\
\hline
Tests E2E web & 0 & 66 \\
\hline
Tests mobiles & 0 & $\sim$30 \\
\hline
Tests de performance & 0 & 3 scripts \\
\hline
Tests de sécurité & 0 & 16 \\
\hline
Pipeline CI/CD avec tests & Non & Oui \\
\hline
Rapports visuels (Allure) & Non & Oui \\
\hline
Couverture code (admin-panel) & 0\% & 57,78\% \\
\hline
\end{tabular}
\end{table}

Au-delà des chiffres, le robot de tests apporte une démarche qualité durable à
l'entreprise. Chaque modification du code est désormais automatiquement validée
par le pipeline CI/CD, réduisant significativement le risque de régressions en
production.

\section*{Compétences acquises}

Ce stage nous a permis de développer des compétences techniques et
méthodologiques variées~:

\begin{itemize}
  \item Maîtrise des frameworks de test modernes (Playwright, JUnit~5, Patrol, k6)
  \item Compréhension approfondie de l'architecture microservices
  \item Pratique du développement piloté par les tests (TDD)
  \item Intégration continue et déploiement continu (CI/CD)
  \item Méthodologie agile et travail en équipe
  \item Rédaction de documentation technique
  \item Analyse de la sécurité applicative (OWASP)
\end{itemize}

\section*{Perspectives et travaux futurs}

Plusieurs axes d'amélioration peuvent être envisagés pour enrichir le robot de
tests~:

\begin{enumerate}
  \item \textbf{Augmentation de la couverture de code} --- Atteindre un objectif
    de 80\% de couverture sur l'ensemble des services, en ciblant
    particulièrement les chemins critiques (paiement, livraison).

  \item \textbf{Tests de contrat évolutifs} --- Implémenter des tests de contrat
    avec Pact pour garantir la compatibilité entre microservices lors des
    mises à jour indépendantes.

  \item \textbf{Ferme de devices mobiles} --- Intégrer une ferme de devices
    (Firebase Test Lab ou BrowserStack) pour exécuter les tests Patrol sur
    de vrais appareils Android et iOS dans le pipeline CI.

  \item \textbf{Tests de régression visuelle} --- Établir les baselines de
    screenshots et activer la comparaison automatique à chaque merge
    request.

  \item \textbf{Monitoring en production} --- Compléter les tests par du
    monitoring applicatif (Grafana, Prometheus) pour détecter les anomalies
    en temps réel.

  \item \textbf{Tests de chaos engineering} --- Introduire des tests de
    résilience simulant des pannes de services (Chaos Monkey) pour valider
    le comportement du système en conditions dégradées.
\end{enumerate}

\vspace{1cm}

En définitive, ce projet de fin d'études nous a offert une expérience
professionnelle riche et formatrice. La conception et le développement d'un robot
de tests automatisés pour un écosystème aussi complexe que Speed Line et Delivery
App nous a confrontés à des défis techniques réels et nous a préparés à notre
future carrière d'ingénieurs en informatique.
